% Cover letter for the Neurocomputing submission.
%
% Compile standalone:  pdflatex cover_letter
% (it does NOT \input any section and does not touch main.tex)
%
% NOTES FOR THE AUTHOR -- read before sending, then delete this block.
%   * Fill in the date and the manuscript ID once Editorial Manager assigns one.
%   * The suggested reviewers are drawn from the recent Neurocomputing papers this
%     manuscript cites, so they are topic-adjacent and demonstrably active in the
%     venue. That is the point of the list, but it also means you must check each
%     one for a conflict of interest (co-authorship, same institution, recent
%     collaboration, advisor/student relation) before entering them. Editorial
%     Manager will ask you to declare any. Drop anyone you are unsure about --
%     a short honest list beats a long risky one.
%   * The letter deliberately does NOT assert "I have no relationship with any of
%     them": that is a factual claim about you that only you can make. If it is true,
%     add it back next to the list; if you are unsure about anyone, remove that name.
%   * Neurocomputing also accepts a short pre-submission inquiry. If you want the
%     scope question settled before spending the submission, send paragraphs 1-2 and
%     the scope paragraph to the Editor-in-Chief first; the rest still stands.

\documentclass[10pt]{article}
\usepackage[margin=0.9in]{geometry}
\usepackage{parskip}
\usepackage{enumitem}
\usepackage{hyperref}
\pagestyle{empty}
\setlist[itemize]{topsep=2pt,itemsep=1pt,parsep=0pt,leftmargin=1.2em}
\linespread{0.98}

\begin{document}

\begin{flushright}
Jiang Tao\\
School of Software Technology, Zhejiang University, Ningbo, China\\
\href{mailto:jiangtao.ml@foxmail.com}{jiangtao.ml@foxmail.com}
\end{flushright}

\vspace{-1em}
\noindent 1 October 2026

\vspace{-0.5em}
\noindent The Editor-in-Chief, \emph{Neurocomputing}

\vspace{-0.5em}
\noindent Dear Editor,

I am submitting \textbf{``FOLoRA: Fisher-Orthogonalized Low-Rank Adaptation for Continual
Learning''} for consideration as a research paper. It is original work, not under
consideration elsewhere, and has a single author.

\textbf{The gap.} Continual learning with pre-trained models is now led by
parameter-efficient methods. Within them, one line --- orthogonalization and subspace
allocation --- decides \emph{where} a new task may write the weights: O-LoRA orthogonally
initializes each task's adapter against earlier tasks' input subspaces, and InfLoRA places
the new subspace in the intersection of its gradient space with the orthogonal complement
of the old tasks'. Both inherit two choices from the classical gradient-projection
literature: the protected subspace comes from a \emph{structural} criterion (an input
covariance, a gradient subspace, an explicit orthogonality constraint) rather than from
the curvature of the past loss, and once chosen, every retained direction is defended
\emph{equally} --- the budget decides how many directions to keep, never how much each one
matters. This manuscript starts from the observation that the second choice is not forced
by anything.

\textbf{What the paper does.} I derive a second-order forgetting bound for sequential
low-rank adaptation, with an explicit higher-order remainder, and show that its leading
term is a Fisher-weighted projection of the adapter displacement onto the past tasks'
eigendirections. That term is exactly what a rank-$k$ regularizer can minimize, which
yields the proposed objective: keep one fixed-budget adapter, estimate a per-layer
importance kernel from the empirical Fisher restricted to it, and penalize displacement
along its top eigendirections in proportion to their curvature. The equal-weight geometry
underlying O-LoRA-style methods is recovered as the special case of equal eigenvalues, so
the contribution generalizes the existing subspace family rather than competing with it.
The method is replay-free, adds no inference cost (the update merges back into the frozen
backbone), and keeps the trainable parameter count constant in the number of tasks,
unlike per-task-adapter designs.

\textbf{Why it is not incremental, against the closest work.} The contribution is
deliberately narrow and therefore checkable. It is not an accuracy claim against EWC: the
paper reports \emph{parity} with the strongest regularized baseline and says so. The
positive results are (i) against the parameter-efficient family it belongs to --- L2P,
CODA-Prompt, InfLoRA, O-LoRA --- where paired differences are significant on both
benchmarks and of one sign across seeds; (ii) an ablation isolating the mechanism:
removing the Fisher-weighted penalty and changing nothing else returns the method to
unprotected sequential fine-tuning, so the entire effect is attributable to the term the
derivation identifies; and (iii) evidence that the protected rank is a memory--accuracy
dial rather than a sensitive hyperparameter. The theory is offered as a statement of
\emph{which} quantity to penalize, not as a numerical predictor --- the manuscript
measures the bound's slack and reports that it is not tight.

\textbf{Scope.} The work is a learning rule: a training objective derived for a
parameter-efficient continual learner and evaluated on standard class-incremental
benchmarks under the frozen-feature nearest-class-mean protocol used by the methods it
compares against. That places it in the journal's interest in learning algorithms and
training objectives, and it engages explicitly with recent \emph{Neurocomputing} work on
task-adaptive and subspace-based continual learning, gradient-based task organization, and
evaluation-protocol artifacts in class-incremental learning, all of which are cited and
discussed in the manuscript.

\textbf{Two places where the paper reports against its own interest} --- the parts I would
most want a referee to check. First, it reports the reproducibility floor it measured
(1.29 accuracy points at a fixed seed and configuration) and then applies it. The
difference against the strongest regularized baseline falls \emph{inside} that band on the
primary benchmark, so that comparison is reported as parity rather than a win, and the
baseline's full regularization curve is printed because the sign of the comparison flips
across it --- including the operating points at which that baseline forgets \emph{less}
than the proposed method. Second, it reports that on one benchmark two reproduced
baselines land at or below the frozen-feature floor, with the mechanism the author
believes is responsible, rather than reporting only the variants that flatter the
comparison. Where a baseline was run on fewer seeds than the proposed method, the table
states each budget rather than implying a shared one. The positive accuracy comparisons
are the ones against the parameter-efficient family, and they are stated as such, not as
a win against full-parameter regularization. Code, protocol and the scripts that
regenerate every table and figure are openly available; the datasets are public.

\textbf{Suggested reviewers.} These researchers have recently published closely related
work in \emph{Neurocomputing}, which is why I suggest them; I would welcome any
alternative the editorial office prefers.
\begin{itemize}
  \item Hongwei Zhao --- task-adaptive LoRA teachers for continual learning.
  \item Lei Pan --- task-adaptive continual learning of vision--language models via
        prototype routing and prompts.
  \item Weizhuo Zhang --- class-incremental learning under classifier bias and calibration.
  \item Junlong Huang --- reducing catastrophic forgetting via gradient-based task
        organization.
  \item Xulong Wang --- continual learning for dissimilar tasks via gradient-norm
        regularization.
\end{itemize}

\textbf{Declarations.} No competing financial interests or personal relationships could
have influenced this work; no external funding; datasets publicly available; code and
reproduction scripts openly available.

\vspace{0.5em}
Thank you for your time and consideration.

\vspace{1em}
Sincerely,\\
Jiang Tao

\end{document}
